\subsection{Functions measurable with respect to a product of two measures}

PROPOSITION 2. --- Let f be a ν-measurable function defined on T×T', with values in a topological space G, and let M be the set of t ∈ T such that the mapping t' ↦ f(t, t') is not μ'-measurable.

\begin{enumerate}
\def\labelenumi{\alph{enumi})}
\item
  If \(f\) is constant on the complement of a \(\nu\) -moderated subset of \(\mathbf{T} \times \mathbf{T}'\) , then \(\mathbf{M}\) is \(\mu\) -negligible.
\item
  If \(\mu'\) is moderated, then M is \(\mu\) -negligible.
\end{enumerate}

The assertion a) follows from Prop. 4b) of §3, No.~2 and the remarks of No.~1. To treat b), note that \(\mu'\) is the sum of a sequence \(\mu_n'\) of bounded measures (§2, No.~3, Prop. 4); \(f\) is measurable with respect to \(\mu \otimes \mu_n' \leqslant \nu\) , and the set M is the union of the sets \(M_{n}\) associated with the measures \(\mu_{n}^{\prime}\) (§2, No.~2, Prop. 2). One is thus reduced to the case that \(\mu^{\prime}\) is bounded, which follows from Prop. 4c) of §3, No.~2.

This statement extends immediately to complex measures (Ch. III, §4, No.~2, Prop. 3).

COROLLARY. --- Let A be a \(\nu\) -measurable subset of \(T \times T'\) , and let M be the set of \(t \in T\) such that the section \(\mathrm{A}(t)\) of A at t is not \(\mu'\) -measurable.

\begin{enumerate}
\def\labelenumi{\alph{enumi})}
\item
  If A is \(\nu\) -moderated, then M is \(\mu\) -negligible.
\item
  If the projection of A on \(T'\) is \(\mu'\) -moderated, then M is locally \(\mu\) -negligible.
\end{enumerate}

The assertion a) follows immediately from Prop. 2. To establish b), denote by B a set, the union of a sequence of \(\mu'\) -integrable open sets in T', that contains the projection of A on T', and denote by \(\mu_1'\) the moderated measure \(\varphi_{\mathrm{B}} \cdot \mu'\) ; since A is measurable with respect to \(\mu \otimes \mu_1' \leqslant \mu \otimes \mu'\) , Prop. 2 implies that A(t) is \(\mu_1'\) -measurable, except for t forming a locally \(\mu\) -negligible set. But since A(t) ⊂ B, to say that A(t) is \(\mu_1'\) -measurable is equivalent to saying that A(t) is \(\mu'\) -measurable (§5, No.~3, Cor. of Prop. 4).

PROPOSITION 3. --- Let f be a mapping of T into a topological space F. If f is \(\mu\) -measurable, then the mapping \((t, t') \mapsto f(t)\) is \(\nu\) -measurable. Conversely, if \(\mu' \neq 0\) , and if this mapping is \(\nu\) -measurable, then the function f is \(\mu\) -measurable.

The first assertion follows from Cor. 1 of Prop. 10 of §6, No.~6. Suppose that \(\mu' \neq 0\) , denote by \(\mu_1'\) a nonzero measure with compact support that is bounded above by \(\mu'\) , by \(\nu_1\) the measure \(\mu \otimes \mu_1'\) , and set \(a = \| \mu_1'\|\) . The projection \(\mathrm{pr}_1\) of \(\mathrm{T} \times \mathrm{T}'\) onto \(\mathrm{T}\) is then \(\nu_1\) -proper, and the image measure \(\mathrm{pr}_1(\nu_1)\) is equal to \(a\mu\) (§6, No.~6, Prop. 10). If \((t, t') \mapsto f(t)\) is \(\nu\) -measurable, then it is also \(\nu_1\) -measurable, therefore \(f\) is measurable with respect to the measure \(a\mu\) (§6, No.~2, Prop. 3), whence the result since \(a \neq 0\) .

The preceding statement extends immediately to complex measures (Ch. III, §4, No.~2, Prop. 3), as do the following corollaries.

COROLLARY 1. --- Let F, F' and G be three topological spaces, and let u be a continuous mapping of F × F' into G. Let f (resp. f') be a function defined on T (resp. T') with values in F (resp. F') and measurable for μ (resp. μ'). Then the function \((t, t') \mapsto u(f(t), f'(t'))\) is measurable for μ ⊗ μ'.

The mappings \((t,t^{\prime})\mapsto f(t)\) , \((t,t^{\prime})\mapsto f^{\prime}(t^{\prime})\) being \(\nu\) -measurable by Prop. 3, this follows from Th. 1 of Ch. IV, §5, No.~3.

COROLLARY 2. --- If A ⊂ T and A' ⊂ T' are measurable (for μ and μ', respectively), then A × A' is measurable for μ ⊗ μ'.

This follows at once from Cor. 1.

COROLLARY 3. --- Consider two positive numerical (resp. complex-valued) functions f defined on T and \(f'\) defined on \(T'\) . If these functions are measurable for \(\mu\) and \(\mu'\) , respectively, then the function

\[
f \otimes f ^ {\prime}: (t, t ^ {\prime}) \mapsto f (t) f ^ {\prime} (t ^ {\prime})
\]

is measurable for \(\mu \otimes \mu'\) .

The case of complex functions, or of finite real functions, is an immediate consequence of Cor. 1. To treat the case of positive numerical functions, for every integer \(n \geqslant 0\) we set \(f_n = \inf(f, n)\) , \(f_n' = \inf(f', n)\) , and we have (with the usual convention \(0 \cdot (+\infty) = 0\) ) \(f \otimes f' = \sup_n(f_n \otimes f_n')\) , whence the result.

PROPOSITION 4. --- Let A be a subset of T. If A is locally \(\mu\) -negligible, then \(A \times T'\) is locally \(\nu\) -negligible. Conversely, if \(A \times T'\) is locally \(\nu\) -negligible and if \(\mu' \neq 0\) , then A is locally \(\mu\) -negligible.

The first assertion follows from Cor. 1 of Prop. 10 of §6, No.~6. To establish the second assertion, let us take up again the notations in the proof of Prop. 3; \(A \times T' = \overline{\text{pr}}_{1}(A)\) is locally negligible for the measure \(\nu_{1}\) , therefore A is locally negligible for \(a\mu\) (§6, No.~2, Cor. of Prop. 2), whence the result since \(a \neq 0\) .

The preceding statement extends at once to the product of two complex measures (Ch. III, §4, No.~2, Prop. 3), as does the following corollary.

COROLLARY. --- If the measure \(\mu\) (resp. \(\mu'\) ) is concentrated on M (resp. M'), then \(\mu \otimes \mu'\) is concentrated on M × M'.

For, \((\mathrm{T} \times \mathrm{T}') - (\mathrm{M} \times \mathrm{M}')\) is the union of the sets \((\mathrm{T} - \mathrm{M}) \times \mathrm{T}'\) and \(\mathrm{T} \times (\mathrm{T}' - \mathrm{M}')\) , which are locally negligible for \(\mu \times \mu'\) by Prop. 4.

